\scp{\tsc{Northwest Seram} (Alune)}{10-02}
Alune,
along with related languages Naka{\Q}ela
and Watui (and sometimes Lisabata),
of the \tsc{Northwest Seram} node is an important witness on
several points, including PMP *-aw > \it{-akw-e}; *b = \it{b},
whereas *mb > \it{p} (devoiced); along with possible trace retentions
of PMP *-uy, *q, and *h.

Alune has many varieties
and much variation, including in sound correspondences.
We have already seen in \srf{10-01} that northern varieties of
Alune have almost universal *d/*z/*R/*l/*j > **l = \it{l} while
western varieties of Alune have a subsequent split of **l > \it{l,
r, d} with complex conditioning environments.
Other differences can be seen in words such as *punti ‘banana’
which becomes \it{uri} (Murnaten,
Riring, Nurue, Buria, Sapolewa varieties), \it{udi} (Hukuanakota), and \it{uʧi} (Kairatu).
Unless otherwise noted, data in this section are from the Laturake
variety of Alune (supplemented occasionally from Riring and Upper Sapolewa).
In examples throughout this section,
\it{-e} (with a variety of onsets) is an extrametrical marker
on many nouns  (not shifting the stress from the root),
and \it{-(n)i} is a genitive enclitic reflecting a part-whole relationship.

Example \qf{10-16} shows the retention of voiced obstruent *b = \it{b}.
In nearly all other \tsc{Ambon-Seram} languages *b undergoes lenition.
The only other languages in which *b does not undergo lenition
are Kayeli (*b = \it{b}),
Wemale (*b > \it{p}),
and Hulung (*b > \it{p}).

\noindent{\wg\st{6pt}\tblr{>{\hspace{-\tabcolsep}}l<{\hspace{-\tabcolsep}}>{\hspace{-\tabcolsep}\enspace}lll}
\lsptoprule
%\tbx{10-16}	&	PMP	&	Alune	&	gloss	\\	\midrule
%\endfirsthead
\tbx{10-16}	&	PMP	&	Alune	& local	gl.	\\	\midrule
%\endhead
	&	*bulan	&	\tbr{b}ulan-e	&	moon	\\
	&	*batu	&	\tbr{b}atu	&	rock	\\
	&	*buaq	&	\tbr{b}ua	&	fruit	\\
	&	*buku	&	\tbr{b}uku-ti	&	bamboo joint	\\
	&	*bubu	&	\tbr{b}u\tbr{b}u	&	fish trap	\\
	&	*baqəRu	&	\tbr{b}elu {\tl} \tbr{b}elu-kwe	&	new	\\
	&	*b\<in\>ahi	&	\tbr{b}ina	&	female	\\
	&	*babaw	&	\tbr{b}a\tbr{b}a-i	&	above, on top	\\
	&	*bətaw	&	\tbr{b}eta-i	&	opposite sex sibling	\\
	&	*babuy	&	\tbr{b}a\tbr{b}u / (apale)	&	pig 	\\
	&	*təbuh	&	te\tbr{b}u	&	sugarcane	\\
	&	*habaRat	&	\tbr{b}alat-e	&	west monsoon	\\
	&	*qabaRa	&	\tbr{b}ala-i	&	arm	\\
	&	*bətəŋ	&	\tbr{h}utun-e	&	foxtail millet (loan? cf. Ambon \it{hotoŋ})	\\
\lspbottomrule
\end{tabular}\mbox{}\\}

By contrast,
the data in example \qf{10-17} show the lenition of *p which
is common in \tsc{Ambon-Seram} languages.
In Alune *p > \it{h,
{\0}}.

\noindent{\wg\st{6pt}\tblr{>{\hspace{-\tabcolsep}}l<{\hspace{-\tabcolsep}}>{\hspace{-\tabcolsep}\enspace}lll}
\lsptoprule
%\tbx{10-17}	&	PMP	&	Alune	&	gloss	\\	\midrule
%\endfirsthead
\tbx{10-17}	&	PMP	&	Alune	&	gloss	\\	\midrule
%\endhead
	&	*pajay	&	\tbr{h}ala, ala	&	rice	\\
	&	*pija	&	ila	&	how many	\\
	&	*pusəj	&	use	&	navel	\\
	&	*punti	&	\tbr{h}uri, uri, uʧi	&	banana	\\
	&	*hapuy	&	au-e	&	fire	\\
	&	*qatəp	&	ate	&	thatch, roof	\\
%	&	*paRih	&	\tcg{(pari)}	&	stingray (Malay; Alune *R > \it{l})	\\
\lspbottomrule
\end{tabular}\mbox{}\\}

Nasal-labial clusters resulting from PMP *ma-p,
*-mp- and *-mb- merge
and are devoiced in Alune to \it{p},
as shown in example \qf{10-18}.
Given the retention of voicing in PMP *b = \it{b},
this merger and devoicing to \it{p} runs
counter to \citeauthor{Blust2009}’s (\citeyear[72]{Blust2009})
claims about “postnasal voicing of stops” in such clusters
in this region.

\noindent{\wg\st{6pt}\tblr{>{\hspace{-\tabcolsep}}l<{\hspace{-\tabcolsep}}>{\hspace{-\tabcolsep}\enspace}lll}
\lsptoprule
%\tbx{10-18}	&	PMP	&	Alune	&	gloss	\\	\midrule
%\endfirsthead
\tbx{10-18}	&	PMP	&	Alune	&	gloss	\\	\midrule
%\endhead
	&	*ma-putiq	&	\tbr{p}uti, \tbr{p}uti-le	&	white	\\
	&	*ma-pənuq	&	\tbr{p}enu, \tbr{p}enu-te	&	full	\\
	&	*umpu	&	u\tbr{p}u-i	&	grandparent	\\
	&	*Rumbia	&	\tbr{p}ia	&	sago	\\
\lspbottomrule
\end{tabular}\mbox{}\\}

One innovation that sets Alune apart
from related languages is
PMP *w > \it{kw},
as in example \qf{10-19}.
An exception to this change
is *qasawa > \it{sau-i} ‘spouse’.
\mbox{}\\

\noindent{\wg\st{6pt}\begin{tabularx}{\linewidth}{>{\hspace{-\tabcolsep}}l<{\hspace{-\tabcolsep}}>{\hspace{-\tabcolsep}\enspace}ll>{\raggedright\arraybackslash}X}
\lsptoprule
%\tbx{10-19}	&	PMP	&	Alune	&	gloss	\\	\midrule
%\endfirsthead
\tbx{10-19}	&	PMP	&	Alune	&	gloss	\\	\midrule
%\endhead
	&	*walu	&	\tbr{kw}alu	&	eight	\\
	&	*siwa	&	si\tbr{kw}a	&	nine	\\
	&	*hawak	&	a\tbr{kw}a-i	&	waist	\\
	&	*bulaw-an	&	bula\tbr{kw}an-e	&	gold	\\
	&	*waRəj	&	\tbr{kw}eli-ti	&	string	\\
	&	*wahiyəR	&	\tbr{kw}el-e	&	water, river	\\
	&	*ka-wanaN	&	ma\tbr{kw}anal-e	&	right (side, hand)	\\
	&	*ka-lawaq	&	\tbr{kw}ala\tbr{kw}a bueti	&	spider	\\
	&	*walay (\tsc{PAn}-F)	&	lusu \tbr{kw}ala-i	&	rib (single rib)	\\
	&	**nawa (PCM)	&	na\tbr{kw}a ai-ni	&	arenga sugar palm tree	\\
	&	**tewa (PCM)	&	te\tbr{kw}a mo	&	dumb (cf. Buru \it{tewa moo} not know)	\\
	&	**niwəR (PCM)	&	ni\tbr{kw}el-e	&	coconut (cf. Kayeli \it{niwel-e} among others)	\\
	&	*ma-huab	&	ma\tbr{kw}a	&	yawn	\\
	&	*huaji	&	\tbr{kw}ali	&	sibling (same sex)	\\
\lspbottomrule
\end{tabularx}\mbox{}\\}

Similar to the retention of *w > \it{kw},
in Alune there is a conditioned retention of PMP *-aw >  \it{akw-e}.
On the limited data available,
it appears that inalienably possessed nouns have *-aw > \it{a},
and others *-aw > \it{-akw-e},
treating the historical *w as a consonant.
\citet[47--50]{Collins1983} sees the presence of the noun marker \it{-e},
or the verbal applicative \it{-ke} as critical for the triggering
of \it{kwe{\#}} in Alune.
Examples of *aw > \it{akwe} are given in \qf{10-20}
and examples of *aw > \it{a} are given in \qf{10-21}.
Sources are as follows T = \citet{TaguchiTaguchi-Grimes1990},
N = \citet{Niggemeyer1951b}, C = \citet{Collins1983},
S = \citet{Stokhof1981b}.

\noindent{\wg\st{6pt}\tblr{>{\hspace{-\tabcolsep}}l<{\hspace{-\tabcolsep}}>{\hspace{-\tabcolsep}\enspace}llll}
\lsptoprule
%\tbx{10-20}	&	PMP	&	Alune	&	gloss	&	source	\\	\midrule
%\endfirsthead
\tbx{10-20}	&	PMP	&	Alune	&	gloss	&	source	\\	\midrule
%\endhead
	&	*labaw	&	malab\tbr{akw}-e	&	rat, mouse	&	(C)	\\
	&		&	ndab\tbr{akw}-e	&	rat, mouse	&	(T)	\\
	&	*kasaw	&	as\tbr{akw}-e	&	rafter	&	(S/C/T)	\\
	&	*kalaw	&	al\tbr{akw}-e	&	hornbill	&	(N/S/C)	\\ \midrule
\tbx{10-21}	&	*bətaw	&	bet\tbr{a}-i	&	opposite sex sibling	&	(T)	\\
	&	*babaw	&	bab\tbr{a}, bab\tbr{a}-i	&	above, on top	&	(N/T)	\\
	&	*ma-linaw	&	merin\tbr{a}	&	calm (Kairatu dialect)	&	(C)	\\
	&	**nakaw	&	man\tbr{a}	&	steal (Kairatu dialect)	&	(C)	\\
\lspbottomrule
\end{tabular}\mbox{}\\}

Closely related Watui,
Naka{\Q}ela, and Lisabata have final \it{ʔwe,
ke, we,} and/or \it{ʔe} in words with final \it{kwe} in Alune.
Examples are given in \trf{10-09},
in which data comes from \citet[49]{Collins1983}.

\begin{table}\tcp{Final \it{-kwe} in Alune compared with related languages}{10-09}
	\begin{threeparttable}\st{3.5pt}
	\begin{tabular}{lll@{ }ll@{ }l}\lsptoprule
*gloss			&	rafter							&	rat, mouse							&	hornbill						&	snake										&	k.o. marsupial\su{†}\\
PMP					&	*kasaw							&	**malabaw								&	*kalaw							&	*nipay									&	**sibaw\\	\midrule
Alune				&	{\hr}ʔasa\tbr{kw}-e	&	{\hr\hr}malaba\tbr{kw}-e&	{\hr}ʔala\tbr{kw}-e	&	{\hr}nia, nia\tbr{kw}-e	&	{\hr\hr}siba\tbr{kw}-e\\
Watui				&	{\hr}wasa\tbr{ʔw}-e	&	{\hr\hr}mab\tbr{oʔ}-e		&	{\hr}aa\tbr{ʔw}-e		&	{\hr}nia\tbr{ʔw}-e			&	{\hr\hr}siba\tbr{ʔw}-e\\
Naka{\Q}ela	&	{\hr}kwasa\tbr{k}-e	&	{\hr\hr}malaha\tbr{w}-e	&											&	{\hr}nia\tbr{w}-e				&	{\hr\hr}siha\tbr{w}-e\\
Lisabata		&	{\hr}asa\tbr{ʔ}-a	&	{\hr\hr}malaha					&	{\hr}ala\tbr{u}-e		&	{\hr}nia								&	\\
		\lspbottomrule
	\end{tabular}
		\begin{tablenotes}\tnsize
			\item[†] {Glossed by \citet[49]{Collins1983} as ‘marsupial sp.’.
								The Alune reflex is glossed by \citet[298]{Niggemeyer1951b} as ‘female cuscus’.
								Compare South Nuaulu \it{siha} ‘k.o.
								cuscus (female)’ and Buru \it{sifa-n} ‘marsupial pouch’.}
		\end{tablenotes}
	\end{threeparttable}
\end{table}

Other sources for final -\it{kwe} in Alune include lexically
specific instances of *q.
There are also two instances of final *k > \it{kwe},
and one instance each of final *h,
*u (via intermediate **uwe),
and *y > \it{kwe}.
Unlike words with final *-aw,
most of these words have alternatives without final \it{kwe}.
Some also have alternative forms with final \it{ke}.
Examples are given in \qf{10-22}.
There are also two possible instances of word medial *q > \it{k,
kwe}; *buqaya > \it{bukaa},\footnote{
The Laturake variety of
Alune \citep{Stokhof1981b} has \it{bukwala} ‘crocodile’ which
appears to have merged the \it{bukaa} form for ‘crocodile’ in
other Alune varieties with \it{kwala} ‘monitor lizard’ (compare
also Buru \it{wela} ‘monitor lizard’).} 
and *maRuqanay > \it{\mbox{makwai},
mokwai} ‘male’.

\noindent{\wg\st{6pt}\begin{tabularx}{\linewidth}{>{\hspace{-\tabcolsep}}l<{\hspace{-\tabcolsep}}>{\hspace{-\tabcolsep}\enspace}l>{\raggedright\arraybackslash}Xl>{\raggedright\arraybackslash}Xl}
\lsptoprule
%\tbx{10-22}	&	PMP	&		&	Alune	&	gloss	&	source	\\	\midrule
%\endfirsthead
\tbx{10-22}	&	PMP	&	*gloss	&	Alune	&	gloss	&	source	\\	\midrule
%\endhead
	&	*daRaq	&	blood	&	lal\tbr{a}, lal\tbr{akw}-e	&	blood, red	&	(N/S/C/T)	\\
	&	*zauq	&	far	&	lau, lau\tbr{kw}-e	&	far, distant	&	(N/S/C/T)	\\
	&	*taRuq	&	store, put away	&	talu, talu\tbr{k(w)}-e	&	fill, order, say	&	(N/S/C)	\\
	&	*anaduq	&	long	&	nanu, nanu\tbr{kw}-e	&	long	&	(N/S/C)	\\
	&	*suksuk	&	pierce, penetrate	&	susu\tbr{k(w)}-e	&	sting, bite (of snake)	&	(N)	\\
	&	Mly. \it{bebek}	&	duck	&	bebe\tbr{kw}-e	&	duck	&	(S)	\\
	&	*tumah	&	clothes louse	&	tuma\tbr{kw}-e	&	louse	&	(S/T)	\\
	&	*baqəRu	&	new	&	belu, belu\tbr{kw}-e	&	new, young	&	(N/S/C/T)	\\
	&	*nipay	&	snake	&	nia, nia\tbr{kw}-e	&	snake	&	(N/S/T)	\\
\lspbottomrule\end{tabularx}\mbox{}\\}

\citet[41ff]{Collins1983} sees that there are also other traces
of PMP *q and *h, which further reveal traces of *-uy.
Following his logic involves several steps.
First, in Alune and related languages, some primary
and secondary vowel sequences *a(C)u > \it{au},
as shown in example \qf{10-23}.

\noindent{\wg\st{6pt}\tblr{>{\hspace{-\tabcolsep}}l<{\hspace{-\tabcolsep}}>{\hspace{-\tabcolsep}\enspace}lll}
\lsptoprule
%\tbx{10-23}	&	PMP	&	Alune	&	gloss	\\	\midrule
%\endfirsthead
\tbx{10-23}	&	PMP	&	Alune	&	gloss	\\	\midrule
%\endhead
	&	*zauq	&	l\tbr{au}, l\tbr{au}kw-e	&	far’	\\
	&	*lahud	&	l\tbr{au}	&	sea, seaward	\\
	&	*pahuq	&	\tbr{au}, \tbr{au}-e	&	mango (Nurue dialect)	\\
\lspbottomrule
\end{tabular}\mbox{}\\}

But Alune
and related languages also show another pattern in secondary
vowel sequences of *a(C)u > \it{oi},
that \citet[41ff]{Collins1983} refers to as a “rounding shift’
which he describes as “Each vowel shifted to the vowel closest
in height but with opposite rounding,
perhaps through misphasing of rounding” (p.41).
This is illustrated in example \qf{10-24}.
\mbox{}\\

\noindent{\wg\st{6pt}\begin{tabularx}{\linewidth}{>{\hspace{-\tabcolsep}}l<{\hspace{-\tabcolsep}}>{\hspace{-\tabcolsep}\enspace}llll>{\raggedright\arraybackslash}X}
\lsptoprule
%\tbx{10-24}	&	Reconstruction	&	*gloss	&	Alune	&	gloss	\\	\midrule
%\endfirsthead
\tbx{10-24}	&	\mc{2}{l}{Reconstruction}	&	*gloss	&	Alune	&	gloss	\\	\midrule
%\endhead
	&	PMP	&	*dahun	&	leaf	&	l\tbr{oi}n	&	leaf	\\
	&	PMP	&	*bahu	&	odour	&	b\tbr{oi}-ni	&	its odour	\\
	&	PCM	&	**taqun{\footnotemark} 	&	full	&	t\tbr{oi}n	&	whole (village)	\\
	&	PCM	&	**saqu	&	sew	&	s\tbr{oi}t nai	&	kind of conical container \\ %made by bending a large leaf into  a cone and then piercing the ends together with short bamboo pegs	\\
	&	PAS	&	**ma(q)u(t)	&	small	&	m\tbr{oi}-te	&	vine with very small leaves	\\
	&	PAS	&	**mba(q)u	&	remove	&	p\tbr{oi}-e	&	blow out, wipe away (mucus)	\\
\lspbottomrule
\end{tabularx}\mbox{}\\}
\footnotetext{
PCM **taqun ‘full’ is probably related to PMP *taqun ‘year’.
\citet[41, 140]{Collins1983} suggests ‘year’ (full,
complete cycle of the sun) became generalised to PCM ‘full,
complete’.}

\citet[42]{Collins1983} recognises conditioning factors that
distinguish the data in example \qf{10-24} from that in \qf{10-23}.
The examples in \qf{10-23} have closed syllables with obstruents.
The examples in \qf{10-24} all have intervocalic consonants that
appear to have been lost,
and either have open syllables,
or end with a continuant.
He says in words from closed syllables with obstruent codas in
\qf{10-23} “no rounding shift occurred”.

Both \citet[139]{Stresemann1927} and \citet[51f]{Collins1983} argue
that Alune also retains traces of PMP *-uy,
such that the full vowel-glide sequence needs to be reconstructed
for the \tsc{Ambon-Seram} subgroup of which Alune is a part.
Their arguments depend on understanding the patterns in \qf{10-24},
and the exceptions in \qf{10-23}.
PMP *hapuy ‘fire’ became Alune \it{au-e} ‘fire’.
It cannot be reanalysed as *[awe] or *[auwe],
because **w > \it{kw} in Alune,
which does not occur in this word for ‘fire’.
PMP *hapuy > \it{au-e} preserves the \it{au} vowel sequence seen
in example \qf{10-23}, but it does not have an obstruent as the coda which conditions that.
It does lose the intervocalic historical consonant *p > {\0}
and if *-uy > \it{u} (as in *babuy > \it{babu} ‘pig’),
then it would end with an open syllable,
and therefore would be expected to follow the pattern of **au
> \it{oi}, which is not found.
So Stresemann and Collins see PMP *hapuy > Alune \it{au-e} ‘fire’
as trace evidence of preserving the full historical vowel-glide
sequence of PMP *-uy as a third pattern different from both \qf{10-23}
and \qf{10-24}.\footnote{
		It appears that Blust (in \citealp{BlustTrussel2020})
		has not understood the significance of \citeauthor{Collins1983}’
		\citeyear[40--52]{Collins1983} data and arguments regarding Alune.
		Thus, for example, Alune \it{boi-ni} ‘its stench’ is assigned
		under PMP *bahuq ‘odor,
		stench’ with an obstruent coda which blocks the vowel shift,
		rather than correctly assigned under the doublet *bahu with an
		open syllable that allows the vowel shift.
		Similarly, by assigning Alune \it{loi-ni} ‘leaf’,
		to PCEMP *daun ‘leaf,
		etc.’, Blust misses Collins’ arguments that the Alune form requires
		the medial consonant of PMP *dahun to have all the pieces together
		to trigger the vowel shift or rounding shift.
		While Collins' discussion of the “rounding shift” is built around
		discussion of “the loss of intervocalic consonants” \citep[41ff]{Collins1983},
		nevertheless on the data available to us,
		the conditioning factors do not seem to require intervocalic *h or *q.
		It appears that defining the codas is sufficient: **au > \it{au}
		/{\gap}C\tsc{[+obstruent]}{\#}; **au > \it{oi} /{\gap}C\tsc{[+continuant]}{\#},
		/{\gap}{\#}.} 
The retention of PMP *b = \it{b},
the devoicing of historical nasal-consonant clusters,
and *-aw > \it{a,
akwe}, make Alune an important witness in eastern Indonesia.
Collins also sees the sporadic retention of traces of *-uy
and *q, *h as significant.
The sporadic retention of traces of *q,
*h is revisited in \srf{13-02-02} in connection with irregular onsets.
